% =============================================================================
%   KIDS WORKBOOK THEME
%   Colourful, friendly palette for children aged 9-12
%   Requires lualatex / xelatex (fontspec)
% =============================================================================

% Emoji support via Segoe UI Emoji (available on Windows)
\newfontfamily\emojifont{Segoe UI Emoji}[Renderer=HarfBuzz]
\newcommand{\emoji}[1]{{\emojifont #1}}

% ---------------------------------------------------------------------------
%   COVER PAGE OVERRIDE
%   A clean, bold kids-workbook cover:
%     - Warm cream background
%     - Orange diagonal lines top-right (geometric accent)
%     - Orange left bar (16 mm)
%     - Large orange title block at the bottom with white title + subtitle
% ---------------------------------------------------------------------------
\renewcommand{\coverpage}[5]{%  #1=title #2=subtitle #3=author #4=date #5=version
  \pagestyle{empty}
  \newgeometry{left=0cm,top=0cm,right=0cm,bottom=0cm}
  
  \definecolor{ModBlue}{HTML}{3B82F6}    % Módulo 1: Azul Eléctrico
  \definecolor{ModGreen}{HTML}{10B981}   % Módulo 2: Verde Esmeralda
  \definecolor{ModYellow}{HTML}{F59E0B}  % Módulo 3: Amarillo / Naranja Cálido
  \definecolor{ModPink}{HTML}{EC4899}    % Módulo 4: Rosa Neón / Magenta
  \definecolor{BgCream}{HTML}{FAF8F2}    % Fondo Crema Claro
  \definecolor{DarkText}{HTML}{1A202C}   % Texto Oscuro Súper Legible

  \begin{tikzpicture}[remember picture, overlay]

      % --- FONDO GENERAL DE LA PÁGINA ---
      \fill[BgCream] (current page.south west) rectangle (current page.north east);

      % --- PATRÓN GEOMÉTRICO ESTILO BAUHAUS A LA IZQUIERDA ---
      \begin{scope}[shift={(current page.north west)}, yshift=-2cm]
          
          % BLOCKS DE PATRÓN REPETITIVO (COLUMNA IZQUIERDA DE MÓDULOS)
          \foreach \i in {0, 1, 2} {
              \begin{scope}[yshift=-\i*8.5cm]
                  % Fila 1: Cuadrados y Círculos Concéntricos (Rosa + Azul)
                  \fill[ModPink, rounded corners=0.4cm] (0, 0) rectangle (3.8, -3.8);
                  \fill[ModBlue] (1.9, -1.9) circle (1.3);
                  \fill[ModYellow] (1.9, -1.9) circle (0.6);

                  % Fila 1 Derecha: Semicírculo Verde
                  \fill[ModGreen] (3.8, 0) rectangle (7.6, -3.8);
                  \fill[BgCream] (3.8, -3.8) -- (3.8, 0) arc (90:0:3.8) -- cycle;

                  % Fila 2: Cuadrados y Círculos (Verde + Amarillo)
                  \fill[ModYellow, rounded corners=0.4cm] (0, -3.8) rectangle (3.8, -7.6);
                  \fill[ModGreen] (1.9, -5.7) circle (1.3);
                  \fill[ModPink] (1.9, -5.7) circle (0.6);

                  % Fila 2 Derecha: Semicírculo Azul con Círculo Amarillo Interno
                  \fill[ModBlue] (3.8, -3.8) rectangle (7.6, -7.6);
                  \fill[BgCream] (3.8, -3.8) -- (7.6, -3.8) arc (0:-90:3.8) -- cycle;
                  \fill[ModYellow] (3.8, -3.8) -- (5.9, -3.8) arc (0:-90:2.1) -- cycle;
              \end{scope}
          }
          
      \end{scope}

      % --- TEXTO SUPERIOR DERECHA (MARCA / SUBTÍTULO) ---
      \node[anchor=north east, align=right, font=\headingfontfamily] at ([xshift=-1.5cm, yshift=-2.5cm]current page.north east) {
          {\fontsize{32}{36}\selectfont \textbf{\color{DarkText} #1}}\\[0.2cm]
          {\fontsize{12}{15}\selectfont \color{DarkText!80} English \& Fun for Kids}\\[0.1cm]
          {\fontsize{10}{13}\selectfont \color{DarkText!60} ENGLISH COURSE #4}
      };

      % --- TÍTULOS VERTICALES GRANDES (ESTILO RECTOR DE LA IMAGEN) ---
      \node[anchor=south east, rotate=90, inner sep=0pt, font=\titlefontfamily] at ([xshift=-1.8cm, yshift=2.5cm]current page.south east) {
          \begin{minipage}{18cm}
              \raggedright
              {\fontsize{48}{54}\selectfont \textbf{\color{ModBlue} WORKBOOK}}\\[0.3cm]
              {\fontsize{42}{48}\selectfont \textbf{\color{ModPink} INTERACTIVE}}
          \end{minipage}
      };

      % --- LÍNEA DE DETALLE Y LEYENDA INFERIOR DERECHA ---
      \draw[ModYellow, line width=2pt] ([xshift=8.5cm, yshift=2.2cm]current page.south west) -- ([xshift=-1.5cm, yshift=2.2cm]current page.south east);

      \node[anchor=south east, align=right, font=\headingfontfamily] at ([xshift=-1.5cm, yshift=1.2cm]current page.south east) {
          {\fontsize{9}{12}\selectfont \color{DarkText!70} \textbf{CONVERSATIONAL A1} $\cdot$ 24 INTERACTIVE SESSIONS}
      };

  \end{tikzpicture}
  \clearpage
  \restoregeometry
}

% ---------------------------------------------------------------------------
%   COLOUR PALETTE
% ---------------------------------------------------------------------------
% We define named colours and then map them to the template semantic tokens.

\usepackage[table]{xcolor}
\usepackage{pagecolor}

% Raw named colours
\definecolor{kidsBlue}   {HTML}{3B82F6}  % vivid sky blue
\definecolor{kidsGreen}  {HTML}{22C55E}  % fresh green
\definecolor{kidsYellow} {HTML}{FCD34D}  % warm yellow
\definecolor{kidsOrange} {HTML}{F97316}  % punchy orange
\definecolor{kidsPink}   {HTML}{EC4899}  % playful pink
\definecolor{kidsPurple} {HTML}{A855F7}  % fun purple
\definecolor{kidsRed}    {HTML}{EF4444}  % alert red
\definecolor{kidsLightBg}{HTML}{FFFFFF}  % white page
\definecolor{kidsCardBg} {HTML}{FFFFFF}  % card white
\definecolor{kidsGray}   {HTML}{6B7280}  % body text
\definecolor{kidsDark}   {HTML}{1F2937}  % headings
\definecolor{kidsGreenEsmeralda}{HTML}{10B981}
\definecolor{kidsWarmYellow}    {HTML}{F59E0B}

% Semantic mapping (overrides the template tokens)
\colorlet{accent}   {kidsBlue}
\colorlet{fg}       {kidsDark}
\colorlet{FG}       {fg}
\colorlet{fgalt}    {kidsGray}
\colorlet{fgacc}    {kidsCardBg}
\colorlet{bg}       {kidsLightBg}
\colorlet{bgalt}    {bg}
\colorlet{bgacc}    {accent}
\colorlet{border}   {kidsDark}
\colorlet{borderalt}{kidsGray}
\colorlet{borderacc}{accent}

\pagecolor{bg}
\color{fg}

% ---------------------------------------------------------------------------
%   FONTS  — friendly rounded look (Windows system fonts)
% ---------------------------------------------------------------------------
% We use Segoe UI (body) and Candara (headings) — both pre-installed on
% Windows and very readable for children.
% (fontspec is already loaded by the base template)

% Body font: Segoe UI — clean, modern, friendly
\setmainfont[
  Ligatures  = TeX,
  BoldFont   = {Segoe UI Bold},
  ItalicFont = {Segoe UI Italic},
  BoldItalicFont = {Segoe UI Bold Italic}
]{Segoe UI}

% Heading font: Bahnschrift — bold, DK-book style
% Bahnschrift is a variable font (single .ttf), so we use AutoFakeBold
% to get a heavy weight for headings and titles.
\renewfontfamily\headingfontfamily[
  Ligatures    = TeX,
  AutoFakeBold = 2
]{Bahnschrift}
\newfontfamily\titlefontfamily[
  Ligatures    = TeX,
  AutoFakeBold = 2.5
]{Bahnschrift}

% ---------------------------------------------------------------------------
%   TYPOGRAPHY RULES
% ---------------------------------------------------------------------------
% No automatic hyphenation — whole words always move to the next line.
\hyphenpenalty=10000
\exhyphenpenalty=10000
\sloppy  % allow slightly looser line-breaking to avoid overflow

% ---------------------------------------------------------------------------
%   EXTRA WORKBOOK-SPECIFIC BOXES (tcolorbox)
% ---------------------------------------------------------------------------
\usepackage[most]{tcolorbox}
\usetikzlibrary{shadows}

% ------ vocabulary card -------
\newtcolorbox{vocabcard}[2][]{%
  enhanced,
  title        = {\bfseries\color{kidsCardBg}#2},
  fonttitle    = \headingfontfamily\large,
  colback      = kidsCardBg,
  colframe     = kidsBlue,
  colbacktitle = kidsBlue,
  coltitle     = kidsCardBg,
  boxrule      = 2pt,
  arc          = 8pt,
  outer arc    = 8pt,
  boxsep       = 8pt,
  left         = 6pt,
  right        = 6pt,
  drop shadow  = {kidsBlue!40},
  #1
}

% ------ activity box ----------
\newtcolorbox{activity}[2][]{%
  enhanced,
  title        = {\bfseries\color{kidsCardBg}\faIcon{pencil-alt}\ #2},
  fonttitle    = \headingfontfamily\normalsize,
  colback      = kidsCardBg,
  colframe     = kidsGreen,
  colbacktitle = kidsGreen,
  coltitle     = kidsCardBg,
  boxrule      = 2pt,
  arc          = 8pt,
  outer arc    = 8pt,
  boxsep       = 8pt,
  left         = 6pt,
  right        = 6pt,
  drop shadow  = {kidsGreen!30},
  #1
}

% ------ song / rhyme box ------
\newtcolorbox{songbox}[2][]{%
  enhanced,
  title        = {\bfseries\color{kidsCardBg}\faIcon{music}\ #2},
  fonttitle    = \headingfontfamily\normalsize,
  colback      = kidsPink!10,
  colframe     = kidsPink,
  colbacktitle = kidsPink,
  coltitle     = kidsCardBg,
  boxrule      = 2pt,
  arc          = 8pt,
  outer arc    = 8pt,
  boxsep       = 8pt,
  left         = 6pt,
  right        = 6pt,
  drop shadow  = {kidsPink!30},
  #1
}

% ------ grammar tip box ------
\newtcolorbox{grammartip}[2][]{%
  enhanced,
  title        = {\bfseries\color{kidsCardBg}\faIcon{lightbulb}\ #2},
  fonttitle    = \headingfontfamily\normalsize,
  colback      = kidsYellow!20,
  colframe     = kidsOrange,
  colbacktitle = kidsOrange,
  coltitle     = kidsCardBg,
  boxrule      = 2pt,
  arc          = 8pt,
  outer arc    = 8pt,
  boxsep       = 8pt,
  left         = 6pt,
  right        = 6pt,
  drop shadow  = {kidsOrange!25},
  #1
}

% ------ reading box ----------
\newtcolorbox{readingbox}[2][]{%
  enhanced,
  title        = {\bfseries\color{kidsCardBg}\faIcon{book-open}\ #2},
  fonttitle    = \headingfontfamily\normalsize,
  colback      = kidsPurple!8,
  colframe     = kidsPurple,
  colbacktitle = kidsPurple,
  coltitle     = kidsCardBg,
  boxrule      = 2pt,
  arc          = 8pt,
  outer arc    = 8pt,
  boxsep       = 8pt,
  left         = 6pt,
  right        = 6pt,
  drop shadow  = {kidsPurple!25},
  #1
}

% ------ answer line helper ---
% \answerline{width}  →  a solid line for writing answers
\newcommand{\answerline}[1][8cm]{%
  \tikz{\draw[kidsGray, line width=0.8pt] (0,0) -- (#1,0);}%
}

% ------ writing lines (blank ruled area) ---
% Solid lines with 22pt spacing so children have room to write comfortably.
\newcommand{\writinglines}[1][5]{%  #1 = number of lines
  \par\vspace{6pt}%
  \foreach \n in {1,...,#1}{%
    \vspace{22pt}%
    \tikz{\draw[kidsGray!60, line width=0.6pt] (0,0) -- (\linewidth,0);}\par%
  }%
  \vspace{6pt}%
}

% ------ vocabulary word pair  (EN | ES) ------
\newcommand{\vocabpair}[2]{%
  \begin{tabular}{@{}p{0.45\linewidth}@{\hspace{0.05\linewidth}}p{0.45\linewidth}@{}}
    {\bfseries\color{kidsBlue}#1} & {\color{kidsGray}#2} \\[2pt]
    \tikz{\draw[kidsBlue!30, line width=0.5pt] (0,0) -- (\linewidth,0);} &
    \tikz{\draw[kidsGray!30, line width=0.5pt] (0,0) -- (\linewidth,0);}
  \end{tabular}
  \vspace{4pt}
}

% ------ star rating (for self-assessment) -----
\newcommand{\selfcheck}[1][3]{%  #1 = filled stars (0-5)
  \foreach \i in {1,...,5}{%
    \ifnum\i>#1
      {\color{kidsGray!40}$\bigstar$}%
    \else
      {\color{kidsYellow}$\bigstar$}%
    \fi
  }%
}

% ------ icon shortcuts (fontawesome5) ----------
\usepackage{fontawesome5}

% ---------------------------------------------------------------------------
%   SECTION TITLE COLOURS  (colourful rotating palette)
% ---------------------------------------------------------------------------
\usepackage{titlesec}

\titleformat{\chapter}[display]
  {\normalfont\titlefontfamily\bfseries\fontsize{28}{32}\selectfont\color{accent}}
  {\chaptertitlename\ \thechapter}{6pt}
  {\bfseries\fontsize{34}{40}\selectfont\color{kidsDark}}

% [right-margin] = 1.5cm forces long chapter titles to wrap before
% touching the right edge — wrapped line keeps the same left alignment.
\titlespacing*{\chapter}{0pt}{10pt}{20pt}[1.5cm]

\titleformat{\section}
  {\normalfont\headingfontfamily\bfseries\fontsize{14}{17}\selectfont\color{accent}}
  {\thesection}{0.8em}{}

\titleformat{\subsection}
  {\normalfont\headingfontfamily\bfseries\fontsize{12}{15}\selectfont\color{kidsDark}}
  {\thesubsection}{0.6em}{}

% ---------------------------------------------------------------------------
%   SIDEBAR PART TITLE HOOK
% ---------------------------------------------------------------------------
\let\kidspart\part
\renewcommand{\part}[1]{%
  \gdef\currentparttitle{\MakeUppercase{#1}}%
  \backgroundtitlevisiblefalse
  \kidspart{#1}%
  \backgroundtitlevisibletrue
}

